\chapter{Preliminarii}
\label{cap:preliminarii}

Acest capitol prezintă noțiunile și tehnologiile pe care se sprijină lucrarea, stadiul actual al
subdomeniului și obiectivele formulate în raport cu acest context.

\section{Tehnologii de bază}

\paragraph{Backend.} Aplicația server este construită în Java~21 cu Spring Boot~3
\cite{springboot}, folosind module pentru web (REST), securitate, validare și acces la date prin
Spring Data JPA / Hibernate. Persistența se face în PostgreSQL~16 \cite{postgresql}, iar schema este
gestionată exclusiv prin migrări Flyway. O disciplină importantă este modul
\code{ddl-auto: validate}: Hibernate nu generează și nu modifică niciodată schema, ci doar o
validează la pornire, ceea ce surprinde din start orice nepotrivire între entități și schema reală.

\paragraph{Frontend.} Interfața este o aplicație de tip \emph{single-page} scrisă în React cu
TypeScript \cite{react}, livrată chiar de către backend de la aceeași origine, ceea ce simplifică
modelul de securitate al cookie-urilor.

\paragraph{Autentificare.} Sistemul folosește \emph{JSON Web Tokens} (JWT) \cite{rfc7519}: un
token de acces de scurtă durată, semnat simetric, transportat în antetul \code{Authorization}, și un
token de reîmprospătare opac, stocat doar ca amprentă criptografică. Parolele sunt stocate ca
amprente generate cu o funcție de derivare a cheilor rezistentă la atacuri de tip dicționar.

\section{Regăsirea informației și căutarea full-text}

\emph{Regăsirea informației} (information retrieval) studiază modul în care se identifică, dintr-o
colecție de documente, pe cele relevante pentru o interogare. Spre deosebire de o potrivire exactă
SQL, căutarea full-text presupune: analiză lexicală (tokenizare, normalizare, eliminarea
diacriticelor), un \emph{index inversat} care leagă termenii de documentele care îi conțin, o funcție
de \emph{relevanță} care ordonează rezultatele și mecanisme de toleranță la erori.

\paragraph{Relevanță și ordonare.} Motoarele moderne folosesc funcția de scor \emph{BM25}
\cite{robertson2009bm25}, o rafinare probabilistică a familiei TF-IDF care echilibrează frecvența
termenilor cu lungimea documentului. Calitatea ordonării se evaluează cu metrici precum
\emph{precision@k} și \emph{nDCG} (normalized discounted cumulative gain) \cite{jarvelin2002ndcg},
care recompensează plasarea documentelor relevante cât mai sus în lista de rezultate.

\paragraph{Toleranța la greșeli.} Căutarea „fuzzy” acceptă potriviri aproximative pe baza distanței
de editare (Levenshtein), esențială pentru greșeli de tastare (de exemplu \emph{benhc}
$\rightarrow$ \emph{bench}). În bazele relaționale, un mecanism echivalent este oferit de extensia
\code{pg\_trgm} (similaritate pe trigrame).

\paragraph{Căutarea în PostgreSQL vs.\ motoare dedicate.} PostgreSQL oferă căutare full-text nativă
prin tipul \code{tsvector}, funcții de ordonare precum \code{ts\_rank\_cd} și indecși GIN
\cite{postgresFTS}. Aceasta este suficientă la scară mică și are avantajul de a păstra totul într-un
singur sistem. Motoarele dedicate, precum OpenSearch \cite{opensearch} (un fork open-source al
Elasticsearch, construit peste biblioteca Lucene), oferă în schimb analizoare configurabile,
sinonime, scoruri BM25, fațete prin agregări și toleranță la greșeli performantă la scară mare.
Alegerea între cele două este o decizie de arhitectură pe care lucrarea o tratează experimental.

\section{Modele de citire derivate și consistență eventuală}

Pentru a separa calea de scriere (tranzacțională, normalizată) de calea de citire (optimizată pentru
interogări), o practică des întâlnită este construirea unui \emph{model de citire derivat} dintr-o
sursă de adevăr unică — o formă relaxată a tiparului CQRS (Command Query Responsibility Segregation)
\cite{kleppmann2017ddia}. Modelul derivat (aici, indexul OpenSearch) este \emph{eventual consistent}
cu sursa de adevăr (PostgreSQL) și poate fi reconstruit oricând. Problema centrală a acestei abordări
este propagarea fiabilă a modificărilor către modelul derivat (problema \emph{dual-write}), discutată
în Capitolul~\ref{cap:contributie}.

\section{Imuabilitatea istoricului}

Ideea că datele istorice trebuie tratate ca fapte imuabile, niciodată suprascrise, este un principiu
de proiectare bine fundamentat în sistemele orientate pe date \cite{helland2015immutability}. În
domeniul de față, un antrenament efectuat este un fapt: el nu trebuie să se modifice atunci când
planul din care a pornit este editat ulterior. Lucrarea implementează acest principiu prin
\emph{instantanee} (snapshots) și chei externe anulabile, detaliate în
Secțiunea~\ref{sec:instantanee}.

\section{Fundamentarea analizorului de programe}

Analiza structurală a programelor pornește de la rezultate din literatura antrenamentului de forță,
în special relația doză–răspuns dintre volumul săptămânal (numărul de serii pe grupă musculară) și
hipertrofie \cite{schoenfeld2017dose}. Aceste rezultate sunt folosite ca \emph{euristici} pentru a
semnala dezechilibre structurale (de exemplu, volum foarte scăzut sau excesiv, lipsa unui tipar de
mișcare), nu ca prescripții individualizate. Analizorul este, prin urmare, determinist, explicabil și
prudent, însoțindu-și mereu rezultatul de un avertisment privind limitele sale.

\section{Obiectivele lucrării în raport cu contextul}

Pornind de la acest context, obiectivele lucrării sunt:
\begin{enumerate}
  \item construirea unei platforme full-stack funcționale, cu integritate puternică a datelor și un
        model de securitate riguros;
  \item proiectarea unui model de citire pentru căutare, derivat și eventual consistent, cu securitate
        care nu poate fi ocolită din parametrii cererii;
  \item evaluarea empirică, pe un corpus scalabil, a momentului în care un motor de căutare dedicat
        devine justificat față de capacitățile native ale PostgreSQL;
  \item demonstrarea corectitudinii prin teste de integrare pe infrastructură reală, nu pe substitute.
\end{enumerate}
